\documentclass[10pt,a4paper]{amsart}
\usepackage[margin=19mm]{geometry}
\usepackage{amsmath,amssymb,mathtools}
\usepackage[scr=rsfs,cal=euler]{mathalfa}
\usepackage{xfrac}
\usepackage[unicode,hidelinks,bookmarks=false]{hyperref}
\newcommand{\ie}{i.e.\ }
\newcommand{\cf}{cf.\ }
\newcommand{\resp}{respectively\ }
\newcommand{\Subsection}{\subsection}
\providecommand{\nobreakdash}{\nobreak\hbox{-}}
\setlength{\emergencystretch}{2em}
\multlinegap0pt
\usepackage{amssymb}
\usepackage{mathtools}
\usepackage[shortlabels]{enumitem}
\usepackage{xcolor}
\usepackage{tikz}
\newtheorem{lemma}{Lemma}
\newcommand{\R}{\mathbb{R}}
\newcommand{\T}{\mathbb{T}}
\title{Structural dichotomy and mass criticality in indirect chemotaxis cascades: fourth-order ellipticity versus Volterra memory}
\author{Jiguang Yu, Louis Shuo Wang}
\date{Source: arXiv:2605.30438v1 (CC BY 4.0)}
\begin{document}
\maketitle

\noindent\emph{Source and licence.} Structural dichotomy and mass criticality in indirect chemotaxis cascades: fourth-order ellipticity versus Volterra memory. Version arXiv:2605.30438v1 (CC BY 4.0), distributed by arXiv under the Creative Commons Attribution 4.0 International licence, http://creativecommons.org/licenses/by/4.0/, as recorded in the arXiv licence metadata for this exact version and shown on the abstract page https://arxiv.org/abs/2605.30438v1. Full licence notice: \texttt{licenses/arxiv-2605.30438-CC-BY-4.0.txt}.\par\smallskip

\noindent\textit{Editorial scope.} Counted unit: the article's lemma lem:heat-Bessel (source0.tex lines 891--902). The article's complete original proof of it is delivered with it (source0.tex lines 904--942, one \texttt{proof} environment of 39 lines). No mathematics is written, repaired or re-derived: the statement and the proof are the article's own lines, in the article's own notation, and no retained line is substituted or re-flowed.  No further source-local context is needed: every cross-reference of every retained line resolves inside the two delivered spans.  Nothing was omitted from any retained span, so the package carries no omission record.  No retained line contains a citation, so the wrapper carries no bibliography and no bibliography entry.  No retained line contains a figure environment, an image-inclusion command or drawing code, so no image is delivered and the package declares no asset dependency.  Editorial formatting only, in addition: an AMS \texttt{amsart} preamble that restates the article's own declarations and macros used by the retained lines, namely the environments \texttt{lemma} on separate counters, together with the standard packages those lines need.  The retained environments print in this document as Lemma 1 in order of appearance, whereas the article prints the same items with the numbering of its own full document.  The complete native source of the article as distributed in the arXiv e-print for this exact version is bundled unchanged as \texttt{sources/201673/source0.tex}.
\begin{lemma}[Heat--Bessel estimate]\label{lem:heat-Bessel}
Let $1<q\le s<\infty$ and $\theta>0$. There exists $C=C(N,q,s)$ such that
\begin{equation}\label{eq:heat-Bessel}
        \|\nabla e^{-\theta A}A^{-1}f\|_{L^{s}(\R^{N})}
        \le C\,e^{-\theta}\theta^{-\beta(q,s)}\|f\|_{L^{q}(\R^{N})},
\end{equation}
with
\begin{equation}\label{eq:beta-qs}
        \beta(q,s):=\Bigl(\tfrac{N}{2}\bigl(\tfrac{1}{q}-\tfrac{1}{s}\bigr)-\tfrac{1}{2}\Bigr)_{+}.
\end{equation}
The analogous estimate (with $e^{-\theta}$ replaced by $1$) holds on $\T^{N}$ and on bounded Neumann domains.
\end{lemma}

\begin{proof}
The Fourier multiplier of $\nabla A^{-1}$ is $i\xi/(1+|\xi|^{2})$, which satisfies the Mikhlin--Hörmander condition $|\xi|^{|\alpha|}|\partial_\xi^{\alpha}(i\xi/(1+|\xi|^{2}))|\le C_\alpha$ for every multi-index $\alpha$. Hence $\nabla A^{-1}$ is bounded on $L^{p}(\R^{N})$ for $1<p<\infty$, and by Bessel-potential / Sobolev embedding it maps $L^{q}\to L^{r}$ boundedly for any $1<q\le r<\infty$ satisfying
\[
        \frac{1}{r}\ge\frac{1}{q}-\frac{1}{N}.
\]
The heat semigroup $\{e^{\theta\Delta}\}_{\theta>0}$ satisfies the standard ultracontractive estimate
\[
        \|e^{\theta\Delta}g\|_{L^{s}(\R^{N})}\le C\theta^{-\frac{N}{2}(\frac{1}{r}-\frac{1}{s})}\|g\|_{L^{r}(\R^{N})},
        \qquad 1\le r\le s\le\infty.
\]
Writing $e^{-\theta A}=e^{-\theta}e^{\theta\Delta}$, the Bessel mass term contributes the exponential factor $e^{-\theta}$, and the heat factor contributes the ultracontractive gain in integrability.

\emph{Case $1/q-1/s\le 1/N$.}
In this range, $\nabla A^{-1}:L^{q}\to L^{s}$ is bounded directly, and $\|e^{-\theta A}h\|_{L^{s}}\le e^{-\theta}\|h\|_{L^{s}}$. Hence
\[
        \|\nabla e^{-\theta A}A^{-1}f\|_{L^{s}}
        =\|e^{-\theta A}\nabla A^{-1}f\|_{L^{s}}
        \le e^{-\theta}\|\nabla A^{-1}f\|_{L^{s}}
        \le Ce^{-\theta}\|f\|_{L^{q}},
\]
which is \eqref{eq:heat-Bessel} with $\beta(q,s)=0$.

\emph{Case $1/q-1/s>1/N$.}
Set $1/r:=1/q-1/N$, so that $\nabla A^{-1}:L^{q}\to L^{r}$ boundedly. Then
\[
        \|\nabla e^{-\theta A}A^{-1}f\|_{L^{s}}
        \le \|e^{-\theta A}\|_{L^{r}\to L^{s}}\,\|\nabla A^{-1}f\|_{L^{r}}
        \le Ce^{-\theta}\theta^{-\frac{N}{2}(\frac{1}{r}-\frac{1}{s})}\|f\|_{L^{q}}.
\]
By the choice of $r$, the exponent is
\[
        \frac{N}{2}\!\left(\frac{1}{r}-\frac{1}{s}\right)
        =\frac{N}{2}\!\left(\frac{1}{q}-\frac{1}{N}-\frac{1}{s}\right)
        =\frac{N}{2}\!\left(\frac{1}{q}-\frac{1}{s}\right)-\frac{1}{2}=\beta(q,s),
\]
yielding \eqref{eq:heat-Bessel}.

The torus and bounded Neumann domain cases proceed identically, using the Mikhlin--Hörmander theorem in those geometries (with the exponential decay $e^{-\theta}$ absent on bounded domains, where it is absorbed into a constant on $[0,T]$).
\end{proof}

\end{document}
